% AY2020 材料力学 〔II〕（転記）
% 出典: 04_材料力学/original/04_材料力学/2020/2020za.pdf
% 要約: 長さ2Lの片持ちはり（固定端A）。AからLの位置Bに曲げモーメントM0,
%       2Lの位置C（自由端）に外力P（下向き）。縦弾性係数E, 断面二次モーメントI。
%       (1) SFD/BMD, (2) BにおけるたわみδB。

\section*{〔II〕}

図2に示すように, 長さ $2L$ の片持ちはりにおいて, A から $L$ の位置 B に曲げモーメント $M_0$,
$2L$ の位置 C に外力 $P$ を加えた. はりの縦弾性係数を $E$, 断面 2 次モーメントを $I$ とする.
次の問い (1) および (2) に答えよ.

\begin{center}
\begin{tikzpicture}[>=Latex]
  % 固定壁（左, A）
  \fill[pattern=north east lines] (-0.4,-0.9) rectangle (0,1.3);
  \draw (0,-0.9) -- (0,1.3);
  % はり（A=0 → C=8, 2L）
  \draw[line width=1.2pt] (0,0.12) -- (8,0.12);
  \draw[line width=1.2pt] (0,-0.12) -- (8,-0.12);
  % B点 曲げモーメント M0（時計回り）
  \draw[->,thick] (3.55,0.55) arc (90:-200:0.55);
  \node[above right] at (4.0,0.45) {$M_0$};
  \fill (4,0) circle (1.4pt);
  % C点 外力 P（下向き）
  \draw[->,very thick] (8,1.3) -- (8,0.18) node[midway,right]{$P$};
  % 節点
  \node[below right] at (0.05,-0.15) {A}; \node[below] at (4,-0.7) {B}; \node[below] at (8,-0.2) {C};
  % 座標 x
  \draw[->] (0,-0.55) -- (1.3,-0.55) node[right]{$x$};
  % 寸法 L, 2L
  \draw[<->] (0,0.85) -- (4,0.85) node[midway,fill=white]{$L$};
  \draw[<->] (0,1.55) -- (8,1.55) node[midway,fill=white]{$2L$};
\end{tikzpicture}

\vspace{1mm}
図2
\end{center}

\begin{enumerate}[label=(\arabic*)]
  \item せん断力図（SFD）および曲げモーメント図（BMD）を示せ.

  \item B におけるたわみ $\delta_B$ を求めよ.
\end{enumerate}
